% GENERATED FILE. Edit canonical lessons, then run npm run generate:books.
% canonical-source-hash: fnv1a64:3d011eef5d3c236a
% canonical-lessons: ZH-C162-yiban, ZH-C162-zongshi, ZH-C162-jingchang, ZH-C162-ouer, ZH-C162-tongchang

\chapter{How Often}
\label{ch:zh-how-often}
\hlchaptermodality{162}

\begin{chapteropening}
\textbf{By the end of this chapter:} \emph{I can say what I do as a rule, every time, regularly, now and then and normally.}

Complete the last lesson of chapter 162: I can say what I do as a rule, every time, regularly, now and then and normally.
\end{chapteropening}

\section[yìbān]{yìbān — generally; as a rule; ordinary}

\label{lesson:ZH-C162-yiban}



\begin{quote}
Before the new one: say the Mandarin for beef, then the Mandarin for chicken meat.
\end{quote}

\subsection*{You'll want to know: yìbān}
\emph{yìbān} — \textquotedblleft{}generally; as a rule; ordinary\textquotedblright{}.

\begin{quote}
\emph{wǒ yìbān qī diǎn qǐchuáng} — As a rule I get up at seven.
\end{quote}

Its characters come later; for now it is a word for the ear.

The tone of \zh{一} changes: yī is said yì before the first tone of bān.

\begin{cousinweb}
\emph{yìbān} is built from \zh{一} 'one' and \emph{bān} 'kind'. The first is inherited from Middle Chinese \emph{'jit} (Guangyun, in Baxter's transcription); the second is listed there as \emph{pan}, glossed 'to carry', and as \emph{ban}, glossed 'joy'. The sense 'kind' is not given, so how the pair came to mean 'as a rule' is not settled.
\end{cousinweb}

\begin{grammarlens}[title={What you've built}]
One more word for times of day.
\end{grammarlens}

\subsection*{Your turn}
\begin{itemize}
\raggedright
  \item \emph{Say it:} \emph{yìbān}
  \item \emph{Say it:} \emph{yìbān}, once more
  \item \emph{Say it:} \emph{wǒ yìbān qī diǎn qǐchuáng} — As a rule I get up at seven.
  \item \emph{Recall:} say the Mandarin for beef, then the Mandarin for chicken meat, then say \emph{yìbān} again
\end{itemize}

\begin{tcolorbox}[breakable,colback=teal!4,colframe=teal!35!black,title={Before you move on}]
What does yìbān mean? (\textquotedblleft{}Generally; as a rule; ordinary\textquotedblright{}.) Say it once more.
\end{tcolorbox}
\section[zǒngshì]{zǒngshì — every single time; always}

\label{lesson:ZH-C162-zongshi}



\begin{quote}
Before the new one: say the Mandarin for chicken meat, then the Mandarin for as a rule.
\end{quote}

\subsection*{You'll want to know: zǒngshì}
\emph{zǒngshì} — \textquotedblleft{}every single time; always\textquotedblright{}.

\begin{quote}
\emph{tā zǒngshì hěn máng} — He's busy all the time.
\end{quote}

Its characters come later; for now it is a word for the ear.

\begin{cousinweb}
\emph{zǒngshì} is built from \emph{zǒng} 'all, gathered together' and \zh{是} 'to be', both inherited from Middle Chinese, \emph{tsuwngX} and \emph{dzyeX} (Guangyun, in Baxter's transcription); how the pair came to mean every time is not settled.
\end{cousinweb}

\begin{grammarlens}[title={What you've built}]
One more word for times of day.
\end{grammarlens}

\subsection*{Your turn}
\begin{itemize}
\raggedright
  \item \emph{Say it:} \emph{zǒngshì}
  \item \emph{Say it:} \emph{zǒngshì}, once more
  \item \emph{Say it:} \emph{tā zǒngshì hěn máng} — He's busy all the time.
  \item \emph{Recall:} say the Mandarin for chicken meat, then the Mandarin for as a rule, then say \emph{zǒngshì} again
\end{itemize}

\begin{tcolorbox}[breakable,colback=teal!4,colframe=teal!35!black,title={Before you move on}]
What does zǒngshì mean? (\textquotedblleft{}Every single time; always\textquotedblright{}.) Say it once more.
\end{tcolorbox}
\section[jīngcháng]{jīngcháng — regularly; often}

\label{lesson:ZH-C162-jingchang}



\begin{quote}
Before the new one: say the Mandarin for as a rule, then the Mandarin for every single time.
\end{quote}

\subsection*{You'll want to know: jīngcháng}
\emph{jīngcháng} — \textquotedblleft{}regularly; often\textquotedblright{}.

\begin{quote}
\emph{wǒ jīngcháng qù Běijīng} — I go to Beijing regularly.
\end{quote}

Its characters come later; for now it is a word for the ear.

\begin{cousinweb}
\emph{jīngcháng} is built from \emph{jīng} and \emph{cháng}, both inherited from Middle Chinese, \emph{keng} and \emph{dzyang} (Guangyun, in Baxter's transcription). The dictionary glosses the first 'constant', using the second as its gloss; the sense 'often' is not given there.
\end{cousinweb}

\begin{grammarlens}[title={What you've built}]
One more word for times of day.
\end{grammarlens}

\subsection*{Your turn}
\begin{itemize}
\raggedright
  \item \emph{Say it:} \emph{jīngcháng}
  \item \emph{Say it:} \emph{jīngcháng}, once more
  \item \emph{Say it:} \emph{wǒ jīngcháng qù Běijīng} — I go to Beijing regularly.
  \item \emph{Recall:} say the Mandarin for as a rule, then the Mandarin for every single time, then say \emph{jīngcháng} again
\end{itemize}

\begin{tcolorbox}[breakable,colback=teal!4,colframe=teal!35!black,title={Before you move on}]
What does jīngcháng mean? (\textquotedblleft{}Regularly; often\textquotedblright{}.) Say it once more.
\end{tcolorbox}
\section[ǒu'ěr]{ǒu'ěr — now and then; occasionally}

\label{lesson:ZH-C162-ouer}



\begin{quote}
Before the new one: say the Mandarin for every single time, then the Mandarin for regularly.
\end{quote}

\subsection*{You'll want to know: ǒu'ěr}
\emph{ǒu'ěr} — \textquotedblleft{}now and then; occasionally\textquotedblright{}.

\begin{quote}
\emph{wǒ ǒu'ěr hē kāfēi} — I drink coffee now and then.
\end{quote}

Its characters come later; for now it is a word for the ear.

Two third tones: ǒu'ěr is said óu'ěr, the first rising.

\begin{cousinweb}
\emph{ǒu'ěr} is built from \emph{ǒu} 'by chance' and the old ending \zh{尔}, both inherited from Middle Chinese (Guangyun, in Baxter's transcription): the first as \emph{nguwH}, glossed 'unexpected', though today's third tone goes with its other reading, \emph{nguwX} 'a pair'; the second as \emph{nyeX}, glossed 'you', here an ending.
\end{cousinweb}

\begin{grammarlens}[title={What you've built}]
One more word for times of day.
\end{grammarlens}

\subsection*{Your turn}
\begin{itemize}
\raggedright
  \item \emph{Say it:} \emph{ǒu'ěr}
  \item \emph{Say it:} \emph{ǒu'ěr}, once more
  \item \emph{Say it:} \emph{wǒ ǒu'ěr hē kāfēi} — I drink coffee now and then.
  \item \emph{Recall:} say the Mandarin for every single time, then the Mandarin for regularly, then say \emph{ǒu'ěr} again
\end{itemize}

\begin{tcolorbox}[breakable,colback=teal!4,colframe=teal!35!black,title={Before you move on}]
What does ǒu'ěr mean? (\textquotedblleft{}Now and then; occasionally\textquotedblright{}.) Say it once more.
\end{tcolorbox}
\section[tōngcháng]{tōngcháng — normally; usually}

\label{lesson:ZH-C162-tongchang}



\begin{quote}
Before the new one: say the Mandarin for regularly, then the Mandarin for now and then.
\end{quote}

\subsection*{You'll want to know: tōngcháng}
\emph{tōngcháng} — \textquotedblleft{}normally; usually\textquotedblright{}.

\begin{quote}
\emph{zhōumò wǒ tōngcháng zài jiā} — At weekends I'm normally at home.
\end{quote}

Its characters come later; for now it is a word for the ear.

\begin{cousinweb}
\emph{tōngcháng} is built from \emph{tōng} 'to go through' and \emph{cháng}, both inherited from Middle Chinese, \emph{thuwng} and \emph{dzyang} (Guangyun, in Baxter's transcription): what holds right through. The sense 'normal' is not given there.
\end{cousinweb}

\begin{grammarlens}[title={What you've built}]
One more word for times of day.
\end{grammarlens}

\subsection*{Your turn}
\begin{itemize}
\raggedright
  \item \emph{Say it:} \emph{tōngcháng}
  \item \emph{Say it:} \emph{tōngcháng}, once more
  \item \emph{Say it:} \emph{zhōumò wǒ tōngcháng zài jiā} — At weekends I'm normally at home.
  \item \emph{Recall:} say the Mandarin for regularly, then the Mandarin for now and then, then say \emph{tōngcháng} again
\end{itemize}

\begin{tcolorbox}[breakable,colback=teal!4,colframe=teal!35!black,title={Before you move on}]
What does tōngcháng mean? (\textquotedblleft{}Normally; usually\textquotedblright{}.) Say it once more.
\end{tcolorbox}
